% use options
%  '[beamer]' for a digital projector
%  '[trans]' for an overhead projector
%  '[handout]' for 4-up printed notes
\documentclass[beamer]{beamer}
\usecolortheme{indiana}

% change talk_preamble if you want to modify the slide theme, colours, and settings for trans and handout modes.
\newcommand{\pathtotrunk}{../}
\input{preamble.tex}
\usepackage{etex}
\usepackage{pgfpages}

%\usepackage{dcpic}
\usepackage{pictex}
\usepackage{pinlabel}
\usepackage{color}

\usepackage{tikz}
\usetikzlibrary{shapes}
\usetikzlibrary{backgrounds}
\usetikzlibrary{decorations.pathreplacing}
\usepackage{tikz-qtree}

% beamer mode
\mode<beamer>{
 \usetheme{CambridgeUS}
}

% transparency mode
\mode<trans>{
 \usetheme{CambridgeUS}
}

% handout mode
\mode<handout>{
 \usetheme{default}
 \setbeamercolor{background canvas}{bg=black!5}
 \pgfpagesuselayout{4 on 1}[letterpaper,landscape,border shrink=2.5mm]
}

\newcommand{\return}[2]{\hyperlink{#1}{\beamerreturnbutton{#2}}}
\newcommand{\goto}[2]{\hyperlink{#1}{\beamergotobutton{#2}}}
\newcommand{\skipto}[2]{\hyperlink{#1}{\beamerskipbutton{#2}}}


% This switches fonts to the Palatino family.
%  \renewcommand{\familydefault}{ppl}

\usepackage{array}

\renewcommand{\pathtodiagrams}{../diagrams/}

%\renewcommand{\bigraph}[1]{{\hspace{-3pt}\begin{array}{c}%
%  \raisebox{-2.5pt}{\includegraphics[height=6mm]{../../graphs/\hashlookup{#1}}}% 
%\end{array}\hspace{-3pt}}}


\usetikzlibrary{matrix}
\usetikzlibrary{arrows,backgrounds,patterns,scopes,external,
    decorations.pathreplacing,
    decorations.pathmorphing
}

%\setbeamertemplate{blocks}[default]

\newlength{\fuzzwidth}
\setlength{\fuzzwidth}{2.5pt}
\newlength{\arrowlength}
\setlength{\arrowlength}{8pt}
\newlength{\arrowwidth}
\setlength{\arrowwidth}{.75pt}
\newlength{\pointrad}
\setlength{\pointrad}{1.5pt}
\newlength{\linewid}
\setlength{\linewid}{1.5pt}
\newlength{\circlerad}
\setlength{\circlerad}{16pt}
\newlength{\smcirclerad}
\setlength{\smcirclerad}{8pt}
\newcommand{\fillcolor}{black!60}
\newcommand{\fuzzcolor}{black!25}
\newcommand{\arrowcolor}{black!25}

\tikzset{
        fuzzright/.style={
        preaction={draw,line width=\fuzzwidth,\fuzzcolor,decorate,decoration={curveto,amplitude=0,raise=-.5*\fuzzwidth}}},
        fuzzleft/.style={
        preaction={draw,line width=\fuzzwidth,\fuzzcolor,decorate,decoration={curveto,amplitude=0,raise=.5*\fuzzwidth}}},
        fuzzrightpre/.style={ %%% Doesn't work
        preaction={draw,line width=2pt,\fuzzcolor,decorate,decoration={curveto,amplitude=0,raise=-1pt,pre=moveto,pre length=12pt}}},
        fuzzleftpre/.style={ %%% Doesn't work
        preaction={draw,line width=2pt,\fuzzcolor,decorate,decoration={curveto,post=moveto,post length=32pt,amplitude=0,raise=1pt}}},        
        outstyle/.style={\arrowcolor, line width=\arrowwidth},
        linestyle/.style={line width=\linewid}
}

\newcommand{\ld}[1]
{{}^*{#1}}
\newcommand{\rd}[1]
{{#1}^*}
\newcommand{\ldd}[1]
{{}^{**}{#1}}
\newcommand{\rdd}[1]
{{#1}^{**}}
\newcommand{\ldddd}[1]
{{}^{****}{#1}}
\newcommand{\rdddd}[1]
{{#1}^{****}}



%\setbeameroption{previous slide on second screen=right}

\author[Noah Snyder]{Noah Snyder (with C. Douglas and C. Schommer-Pries)}
\title{What is a homotopy coherent $SO(3)$ action?}
\date{MSRI: Tensor categories and TQFTs}
\institute{Indiana University}

\titlegraphic{\includegraphics[width=3cm]{iu_trident_web_crimson}
}

%\usepackage{multimedia}


\begin{document}

\frame{\titlepage}

\begin{frame}
       \frametitle{Outline}
       \tableofcontents
\end{frame}

\beamertemplatetransparentcovered 

\beamersetuncovermixins 
{\opaqueness<1->{60}} 
{} 


\AtBeginSection[]
{
   \begin{frame}<beamer>
       \frametitle{Outline}
       \tableofcontents[currentsection]
   \end{frame}
}

\section{Introduction}

\begin{frame}{Local Topological Quantum Field Theories}
\begin{block}{TQFTs informally}
An $n$-dimensional TQFT is an invariant of $n$-manifolds which can be computed via cutting and pasting.
\end{block}

\begin{block}{Local TFTs informally}
\begin{itemize}
\item Better to cut up along lower dimensional pieces
\item Best of all cut down to points.
\end{itemize}
\end{block}

\begin{block}{Formal Definition}
An $n$-dimensional fully local topological field theory with values in a symmetric monoidal (not necessarily discrete) $n$-category  $\cC$ is a symmetric monoidal functor:
$$\cF: \mathrm{Bord}_n \rightarrow \cC.$$
\end{block}
\end{frame}

\begin{frame}{Topological structures}

\begin{block}{Flavors of TFT}
Different topological structures on bordisms give different TFTs:
\begin{itemize}
\item oriented
\item spin
\item framed (choice of trivialization of the tangent bundle)
\item etc.
\end{itemize}
\end{block}

\begin{block}{Lower dimensions}
In order to glue structures, we need to pick a structure on small collars of boundaries.

\vspace{.1in}
For example, a $3$-framing on a $1$-manifold is a trivialization of $TM \oplus \mathbb{R}^2$.
\end{block}

\end{frame}

\begin{frame}{Baez-Dolan Cobordism Hypothesis}

\begin{block}{Dualizable objects}
A symmetric monoidal $n$-category is fully dualizable if every object has a dual and ever $k$-morphism for $1 \leq k < n$ has a left adjoint and a right adjoint.  
\newline

Every symmetric monoidal $n$-category has a maximal fully dualizable subcategory $\cC^{fd}$.
\end{block}

\begin{theorem}[Lurie(-Hopkins) 09]
$$\mathrm{TFT}^{fr}(\cC) \tilde{\rightarrow} \cC^{fd}$$ as \emph{spaces} via $$\cF \mapsto \cF(\mathrm{pt}_+).$$
\end{theorem}
\end{frame}

%\begin{frame}{Where does the cobordism hypothesis come from?}
%\begin{block}{Why is $\mathrm{Bord}_n$ fully dualizable?}
%\begin{itemize}
%\item The dual of the positively framed point is the negatively framed point.
%\item The evaluation and coevaluation are given by $1$-handles.
%\item The adjoints of the $1$-handles are other $1$-handles.
%\item The units and counits of these adjunctions are given by $2$-handles.
%\end{itemize}
%\end{block}
%\begin{block}{Why framed?}
%The framing keeps track of the difference between left and right adjoints.  $\mathrm{Bord}_n^{fr}$ is the {\em universal} fully dualizable category.  Proof uses Igusa's version of generalized Morse theory.
%\end{block}
%\end{frame}

\begin{frame}{$SO(n)$ action}
\begin{block}{$SO(n)$ action on framed TFTs}
$SO(n)$ acts on the space of framed TFTs via changing the framing.
\end{block}

\begin{block}{$SO(n)$ action on fully dualizable objects}
By the cobordism hypothesis there's a ``homotopy coherent action" of $SO(n)$ on the \emph{space} of fully dualizable objects.
\end{block}

\begin{block}{Warnings:}
\begin{itemize}
\item By ``space" we mean keep only the invertible morphisms.  This is an $\infty$-groupoid, i.e. a ``space" by Grothendieck's homotopy hypothesis.
\item Since we've translated across a homotopy equivalence, this action is only defined homotopy coherently.
\end{itemize}
\end{block}
\end{frame}


\begin{frame}{Main question}

\begin{block}{Main question}
What does it mean \emph{explicitly} to give a homotopy coherent action on a $SO(3)$ (discrete) $3$-category?
\end{block}

\begin{block}{Why we care}
\begin{itemize}
\item Use TFTs to give applications from topology to algebra.
\item Fully understand all consequences of the cobordism hypothesis in simple models.
\item Important if you want to study oriented TFTs instead of framed.
\end{itemize}
\end{block}

\begin{theorem}[Lurie]
The homotopy fixed points of the $SO(n)$ action on $\mathrm{Bord}_n$ are exactly the oriented local TFTs.
\end{theorem}

\end{frame}



\begin{frame}{Main example}

\begin{theorem}[Haugseng, Scheimbauer--Johnson-Freyd]
There's a symmetric monoidal ``Morita" $3$-category $\cT\cC$ of tensor categories, bimodule categories, bimodule maps, and bimodule natural transformations.
\end{theorem}

\begin{theorem}[DSPS 13]
Fusion categories (of non-zero global dimension) are fully dualizable objects in $\cT\cC$.
\end{theorem}

\begin{conj}[DSPS 13]
Non-semisimple finite tensor categories are \emph{not} fully dualizable, but there's still an $SO(3)$-action.
\end{conj}
\end{frame}

\begin{frame}{Other examples}
\begin{example}
$SO(2)$-action on separable algebras (Davidovich over $\mathbb{C}$, Patrick Chu for all fields).
\end{example}

\begin{example}
$SO(n)$-action on $E_n$-algebras (Gwilliam--Scheimbauer).  \\ Comes from a natural action of $SO(n)$ on the $E_n$-operad
\end{example}

\begin{example}
$SO(3)$-action on cp-rigid braided tensor categories (Brochier-Jordan-S.)  Warning: this is a $(4,3)$-category.
\end{example}
\end{frame}

\begin{frame}{Outline:}
\begin{block}{Algebraic topology of the $4$-truncation of $BSO(3)$.}
Action is $BG \rightarrow B \mathrm{Aut}(X)$.  The RHS is a $4$-type.  Want to give a ``minimal" CW description of the $4$-truncation of $BSO(3)$.
\end{block}

\begin{block}{Bordism model of the braided monoidal $2$-groupoid coming from $BSO(3)$.}
On the categorical side, this is a functor of braided monoidal $2$-categories.  Describe this ``quantum topologically" using a Pontryagin-Thom bordism model.
\end{block}

\begin{block}{Expanded description of $SO(3)$ action on a discrete $3$-category}
Use ``Morse decompositions" in the bordism model to give a fully explicit version of the action on a $3$-category (suppressing associators).
\end{block}

\begin{block}{Homotopy fixed points.}
\end{block}
\end{frame}

\section{Algebraic topology}

\begin{frame}{What is a homotopy coherent action?}
\begin{definition}
A homotopy coherent action of $G$ on $X$ is a map $BG \rightarrow BAut(X)$.
\end{definition}

\begin{example}
A homotopy coherent action of $SO(2)$ is a map $\mathbb{C}P^\infty \rightarrow BAut(X)$.  $\mathbb{C}P^\infty$ has a CW description with a cell in every even degree.  Giving a map out of $BSO(2)$ is giving the values of each of these cells in $BAut(X)$.
\end{example}

\begin{block}{What about $SO(3)$?}
$BSO(3)$ is complicated/difficult (but see Ayala-Francis).  We can simplify the question using that our category is discrete.
\end{block}
\end{frame}

\begin{frame}{Truncation and actions}
\begin{block}{Truncation}
If $X$ is a homotopy $n$-type (i.e. $\pi_k(X)$ vanishes for $k > n$).  Maps to an $n$-type only depend on the $n$-truncation of the source.  In a CW complex you only need to go to $n+1$ cells to describe the $n$-type.
\end{block}

\begin{block}{Truncations and actions}
Under the homotopy hypothesis, $n$-groupoids correspond to homtopy $n$-types.  If $X$ is an $n$-type,  $Aut(X)$ is an $n$-type, and $BAut(X)$ is an $n+1$-type.  
\end{block}

\begin{example}
For the $SO(2)$-action on algebras we only need to understand the $3$-type of $BSO(2)$, so we only need the first two cells.
\end{example}

\end{frame}

\begin{frame}{Main topological result}
\begin{theorem}[DSPS]
$BSO(3)$ has the same $4$-type as the CW complex with:
\begin{itemize}
\item One $2$-cell,
\item One $3$-cell attached to $S^2$ via the degree $2$ map (yielding $\Sigma \mathbb{R}P^2$),
\item One $4$-cell attached to $S^2$ via the Hopf map (yielding $\mathbb{C}P^2$),
\item One $5$-cell attached to $\Sigma \mathbb{R}P^2$ via the generator of $\pi_4(\Sigma \mathbb{R}P^2) = \mathbb{Z}/4\mathbb{Z}$.
\end{itemize}
Furthermore, any other such CW complex must have at least as many cells.
\end{theorem}

\begin{block}{Homotopy groups of $\Sigma \mathbb{R}P^2$}
Homotopy groups of $\Sigma^k \mathbb{R}P^2$ are described in Jie Wu's monograph ``Homotopy theory of the suspensions of the projective plane."
\end{block}
\end{frame}

\begin{frame}{Brief summary of argument}
\begin{definition}
$BOrp(3)$ is the homotopy fiber of $BSO(3) \ra BSO \xra{p_1} K(\ZZ,4)$.  
\end{definition}

\begin{theorem}
$BOrp(3)$ has the same $4$-type as the CW complex with
\begin{itemize}
\item One $2$-cell
\item One $3$-cell attached to $S^2$ via the degree $2$ map (yielding $\Sigma \mathbb{R}P^2$)
\item One $5$-cell attached to $\Sigma \mathbb{R}P^2$ via the generator of $\pi_4(\Sigma \mathbb{R}P^2) = \mathbb{Z}/4\mathbb{Z}$.
\end{itemize}
\end{theorem}
\begin{proof}
A cohomology calculation using Serre spectral sequences shows that the map $\Sigma \mathbb{R}P^2 \rightarrow BSO(3) \rightarrow BOrp(3)$, is an iso on $H_{\leq 4}$ and thus by Hurewicz on $\pi_{\leq 3}$.  The $5$-cell then assures that the map is an iso on $\pi_4$.
\end{proof}
\end{frame}

\begin{frame}{More brief summary}
In addition to the map $BOrp(3) \rightarrow BSO(3)$ using the $2$, $3$, and $5$ cells, we also have a map $BSO(2) \rightarrow BSO(3)$ using the $2$ and $4$ cells.  So we need to prove.

\begin{theorem}
The following square is a $4$-pushout.
\begin{equation*}
\xymatrix{
S^2 \ar[r] \ar[d] & BOrp(3) \ar[d] \\
BSO(2) \ar[r] & BSO(3)
}
\end{equation*}
\end{theorem}

This is again proved by a cohomology/Serre spectral sequence argument.  Let $P$ be the pushout, we check that $H_{\leq 5}(BSO(3),P) = 0$, thus $\pi_{\leq 5}(BSO(3), P) = 0$, and the map $P \rightarrow BSO(3)$ is an iso on $\pi_{\leq 4}$.
\end{frame}


\begin{frame}{How useful is this?}

\begin{corollary}
Description of the space of maps from $BSO(3)$ to any $4$-type.
\end{corollary}

\begin{block}{Is this useful?}
We want to understand actions on a $3$-groupoid $X$.  To apply the above literally, we'd need to study $BAut(|X|)$ and then translate back into $3$-groupoids.
\end{block}

\begin{block}{Why there's hope}
If $X$ is a $3$-groupoid, there's a $4$-groupoid whose only point is $X$, whose $1$-morphisms are autofunctors, etc.  This $4$-groupoid is equivalent to the fundamental $4$-groupoid of $BAut(X)$.
\end{block}

\begin{block}{HoTT aside (Buchholtz--van Doorn--Rijke)}
In Homotopy Type Theory $BAut(X)$ is described directly as $\Sigma_{T: \cU} | T = X |$.
\end{block}

\end{frame}

\section{10 minute break for questions}

\begin{frame}{Questions?}
Please ask questions over Zoom.
\end{frame}

\section{Bordism models}

\begin{frame}{Goal of this section}
\begin{block}{Goal}
Describe a $4$-groupoid by generators and relations whose fundamental $4$-groupoid is equivalent to $BSO(3)$.
\end{block}

\begin{block}{But I'm afraid of $4$-categories}
Since $BSO(3)$ is $2$-connected, we can rephrase everything in terms of the braided monoidal $2$-category $\Omega^2BSO(3)$.  Braided monoidal $2$-categories are less scary.
\end{block}

\begin{block}{Is this rigorous?}
Most of this can be made rigorous in existing models of higher categories, with some gaps to be filled later.  In this paper, instead of working in a model we will work model independently but with clear statements about what our requirements are on the model.
\end{block}

\end{frame}




\end{document}




Recall we have a map from $\Sigma \mathbb{R}P^2 = M(\ZZ/2, 2) \rightarrow BSO(3)$.

\begin{block}{Homotopy groups of $BSO(3)$}
Because $SO(3) = \mathbb{R}P^3$ we have the homotopy groups of $BSO(3)$:
$\pi_2 = \mathbb{Z}/2$, $\pi_3 = 0$, $\pi_4 = \mathbb{Z}$.
\end{block}